\section{What the corpus supports for each original question}
\label{app:scope}
The project began with eight research questions and three cross-cutting ones. Table~\ref{tab:scope} states, for each, what these twenty games establish and what they leave open. An open entry means the design cannot measure the construct. It is not evidence that an ability is absent.

\begin{longtable}{@{}L{.20\linewidth}L{.40\linewidth}L{.34\linewidth}@{}}
\caption{Supported scope for the eight questions and for A to C.}\label{tab:scope}\\
\toprule Question & Established in this corpus & Not established\\\midrule\endfirsthead
\toprule Question & Established in this corpus & Not established\\\midrule\endhead
1. Strategic reasoning & Refusals and prices argued from a rival's building budget; first financing actions set against the asset bound for all 188 episodes, with the nine departures decomposed & Optimal play, regret, or a strategic score\\
2. Opponent modeling & Stated figures for a counterparty's cash, costs and liquidation value, some exact (G20) and some wrong (G2) & Calibration of beliefs; transfer to unfamiliar opponents\\
3. Negotiation and deception & All 920 responses by form and completion; two countering and two rejecting response styles; a permitted bluff, an omission, a correction, an appeal and a withdrawn assurance & Whether any tactic persuaded; prevalence of deception; welfare or fairness\\
4. Planning and memory & Stated dependencies and ordered routes before the menu offers the resource; a named selection of assets carried out exactly (G13); what notes each agent carried & A general planning ability; any benefit or harm from memory\\
5. Adaptation and learning & An inventory revised inside one response (G5); offers revised within threads & Learning from feedback across games or turns\\
6. Economic judgment and risk & Exact asset bounds in 43 of 48 unpayable first assessments; five listed-only totals; mortgage-first financing that kept buildings in 28 of 30 cases & Least-cost liquidation, reserve policy, or risk-adjusted value\\
7. Collective behavior & One public appeal to deny a rival cash, with the speaker's own refusal and another player's contrary offer & A coalition, collusion, or its effects\\
8. Evaluation validity & That payment here follows from not stopping; undisclosed differences between published and implemented rules; unequal channels across systems & Performance under a corrected interface or disclosed rules\\\midrule
A. Strategic language versus action & Agreement between first action and asset bound; six Grok responses whose text and call disagree; two false budgets acted on & A causal relation between a written assessment and the action\\
B. General versus multi-agent ability & Nothing. The roster has four unequal configurations and no external capability measure was matched to them & Any relation between general ability and play\\
C. Skill versus luck & Wins by system (8, 7, 3, 2) and by seat (7, 7, 3, 3) & Shares of skill, seat and luck in winning\\\bottomrule
\end{longtable}

\section{Why continued financing pays}
\label{app:proof}
This appendix gives the argument of \S\ref{sec:task} in full. It reads the engine source at the two commits used for the games, which are identical for the files involved, and the holdings recorded at the first menu of each payable debt. No policy, engine step, replay or search was run.

\textbf{Claim.} In each of the 131 payable entry states, every uninterrupted sequence of valid financing actions that each raise cash ends with the current debt paid, provided the player does not declare bankruptcy.

\textbf{What the claim excludes.} Invalid calls, empty sale plans, limits on calls or time, provider failures, a later debt in the same payment queue, and any state outside the 131. It is not a statement that an arbitrary agent would pay, and it says nothing about cost.

\textbf{Source facts.} File \path{python/packages/engine/src/monopoly_engine/engine.py}.
\begin{enumerate}
\item \path{_resolve_pending_payment}: after each financing action, if cash covers the amount the engine pays. Otherwise it computes the available mortgages and sales. If neither list has an entry it ends the episode in bankruptcy. If either has an entry it issues another financing menu.
\item \path{_mortgageable_space_indices}: an owned deed can be mortgaged if it is unmortgaged, carries no building, and no deed of its group carries a building.
\item \path{_sellable_building_space_indices}: in a fully owned group, the deeds at the group's highest building level are listed for sale. A hotel is listed only if the bank holds at least four houses.
\item \path{_validate_sell_plan}: a sale plan may bundle items. A house sale returns the house to the bank and pays half its cost. A hotel sale turns the hotel into four houses taken from the bank and pays half of one building cost. The plan is valid if the bank's house stock stays nonnegative at the end and each touched group ends with levels within one of each other.
\end{enumerate}
Financing never adds a building and never lifts a mortgage.

\textbf{Invariant and termination.} Let $Q$ be cash, plus half the cost of every standing building level with a hotel counted as five, plus the mortgage value of every unmortgaged deed. A mortgage moves its value from the third term to the first. A house sale moves half a house cost from the second to the first, and so does a hotel sale, which reduces five levels to four. $Q$ is therefore constant. Each action that raises cash removes at least one building level or one unmortgaged deed, so the sequence is finite.

\textbf{No unpaid dead end without a hotel.} Suppose the sequence stops unpaid because nothing is offered. If no building stands, every unmortgaged deed is mortgageable, so none remains and cash equals $Q\ge D$, a contradiction. If a building stands in a fully owned, evenly built group without a hotel, a highest house is listed. So an unpaid dead end needs a group whose highest level is a hotel while the bank holds fewer than four houses.

\textbf{Coverage of the 131 entries.} At entry, every built group is fully owned and its levels differ by at most one, and $Q\ge D$, in all 131. Four sufficient conditions, applied in order, cover them (Table~\ref{tab:coverage}).

\begin{table}[H]\centering\small
\begin{tabularx}{\linewidth}{@{}L{.36\linewidth}rY@{}}\toprule
Condition at entry & Debts & Why no unpaid dead end can occur\\\midrule
\rows{gen_coverage_rows}
\bottomrule
\end{tabularx}
\caption{Sufficient conditions and the number of payable debts each is the first to cover. All sixteen debts that need released collateral fall under the first (fifteen) or third (one, G20 d327) condition.}\label{tab:coverage}
\end{table}

For the third condition, let $T$ be bank houses minus four times the hotels held. A house sale raises $T$ and a hotel sale leaves it unchanged, so $T\ge 0$ throughout, and a dead end would need a hotel with fewer than four bank houses, where $T<0$.

For the fourth condition, let $R$ be bank houses plus houses standing in groups that hold no hotel. A house sale in a hotel-free group leaves $R$ unchanged. A house sale elsewhere, possible only inside a bundle, raises it. Selling the only hotel of a group takes four houses from the bank and adds at least four to the hotel-free count, so $R$ does not fall. A dead end with a hotel standing has no listed house, so no hotel-free group holds one, and fewer than four bank houses, so $R<4$. The two entries concerned start with $R=8$ (Table~\ref{tab:lastclass}).

\begin{table}[H]\centering\small
\begin{tabular}{@{}lrrr@{}}\toprule
Entry & Bank houses & Houses in hotel-free groups & $R$\\\midrule
\rows{gen_lastclass_rows}
\bottomrule
\end{tabular}
\caption{The two entries under the fourth condition. Both are Gemini in G3 with one hotel in light blue, one in dark blue, and houses on red.}\label{tab:lastclass}
\end{table}

\textbf{Provenance.} All four authors of the earlier reviews derived the result for the sixteen necessary-collateral entries from the source. The first three conditions were stated and counted by this draft's author on 14 November. The fourth is Astra's argument from that day. This draft re-derived it, extended it to bundled plans as above, and recomputed its premises from the raw records. No one has executed a financing policy on these states.

\section{The payable partition}
\label{app:middle}
The 53 debts in the middle row of Table~\ref{tab:partition} have $C+M_0<D\le U_0$. The bound shows that sales were needed. That released mortgages were not needed rests on three kinds of evidence. Forty-seven were paid by routes that used no mortgage absent from the first menu. One, Grok's G12 d249, was declared, and the stored audit paid it by selling five pink houses, leaving \$43. Five were paid by routes that did use a released mortgage. For each of those five, selling buildings alone, with the listed mortgages, covers the debt, and the bank's stock permits it (Table~\ref{tab:middle}).

\begin{table}[H]\centering\footnotesize
\begin{tabularx}{\linewidth}{@{}llrrrrrY@{}}\toprule
Entry & System & $D$ & $C$ & $M_0$ & $S$ & Bank houses & Buildings at entry\\\midrule
\rows{gen_middle_rows}
\bottomrule
\end{tabularx}
\caption{The five middle-row payments that used a released mortgage. In each, $C+M_0+S\ge D$.}\label{tab:middle}
\end{table}

G3 d228 and d356 hold houses only, and every house can be sold in turn. G3 d482 holds one hotel with seventeen bank houses, and G18 d428 three hotels with twenty-nine, so every exchange is affordable. G3 d461 holds two hotels with six bank houses. Exchanging one hotel leaves two bank houses, selling that group's houses restores the stock, and the second exchange follows. These five routes are arguments from recorded holdings and the sale rules. They were not executed.

\section{Evidence tiers for the sixteen first responses}
\label{app:coding}
The tiers are defined over every available channel of the first response, and the first definition that applies is used.
\begin{description}[leftmargin=1.4em,itemsep=1pt]
\item[Explicit dependency.] A channel says that selling or clearing buildings enables, unlocks, allows or must precede mortgaging deeds of that group.
\item[Ordered route.] No such statement, but a channel gives a sequence in which a building sale comes before a named mortgage of a deed in the same group.
\item[Inclusive budget.] Neither of the above, but a stated total or itemization includes the mortgage value of deeds in a built group.
\item[Tentative resource.] The released mortgages are raised as a possibility and not counted.
\item[Omission.] No available channel mentions the mortgage value of deeds in a built group, and the stated maximum excludes it.
\end{description}
One LLM author assigned the tiers with the stored fields in hand. Two other LLM authors arrived at the same five-way split independently in the 14 November review. That is agreement among model readers of the same records. It is not human validation. The build checks that each evidence span occurs verbatim in the stored channel. Table~\ref{tab:coding} adds the arithmetic and the executed route for each case.

{\footnotesize\setlength{\tabcolsep}{4pt}
\begin{longtable}{@{}L{.072\linewidth}L{.052\linewidth}rrrL{.08\linewidth}L{.285\linewidth}L{.185\linewidth}L{.07\linewidth}@{}}
\caption{The sixteen first responses: sums, tier, figures and route. $U_0$ is the listed-only total and $U_*$ the full bound.}\label{tab:coding}\\
\toprule Entry & System & $D$ & $U_0$ & $U_*$ & Tier & Figures and other statements & Route taken & Outcome\\\midrule\endfirsthead
\toprule Entry & System & $D$ & $U_0$ & $U_*$ & Tier & Figures and other statements & Route taken & Outcome\\\midrule\endhead
\rows{gen_sixteen_coding_rows}
\bottomrule
\end{longtable}}

\section{All sixteen first responses, both channels}
\label{app:texts}
Each entry prints the complete authored note and the complete provider reasoning field of the first response in the episode. For Opus the first attempt was refused with no text, and the fields come from the action-only retry, which has no note. Text is reproduced character for character. Paragraph breaks in provider fields are kept. The machine-readable file \path{supplement/sixteen_first_texts.json} adds the public message, the assistant content, field hashes, source lines and the route.

\fieldhead{G1 d476}{Grok, turn 212. Debt \$200; cash \$51; listed-only total \$126; full bound \$526. Applied first action: sell houses or hotel. Paid with \$56 left after 2 actions. Tier: explicit dependency.}
\fieldtext{Authored note (applied tool argument, standard response)}{Short \$149 on RR \$200. Only legal raise is Marvin 1h (\$75) \(\rightarrow\) cash \textasciitilde{}\$126, still \textasciitilde{}\$74 short. Must sell before yellows can be mortgaged. Next prompt: mortgage one yellow (Atlantic/Ventnor \$130 or Marvin \$140) to clear the debt. Do not bankrupt — 1v1, p3 inherits the win. After pay, cash will be thin; no rebuild, no unmortgage. Jail cards stay.}
\fieldtext{Provider field (reasoning summary)}{I need to decide between selling a house on Marvin Gardens or declaring bankruptcy.\par Selling one house on Marvin Gardens would yield \$75. That covers part of the \$149 shortfall, but I still need roughly \$74 more to pay the rent.\par Selling the house on Marvin Gardens would bring my cash to \$126. That still leaves a \$74 shortfall. I must sell this building first because mortgaging a property with buildings in a monopoly requires even distribution of houses. After that, I can mortgage Atlantic, Ventnor, and Marvin Gardens.}
\fieldtext{Assistant content (returned beside the call, not applied)}{Selling the Marvin house is the only way to unlock the yellow mortgages and cover this rent. Bankruptcy hands the game to p3.}

\fieldhead{G2 d216}{Sol, turn 73. Debt \$450; cash \$133; listed-only total \$283; full bound \$683. Applied first action: sell houses or hotel. Paid with \$103 left after 3 actions. Tier: explicit dependency.}
\fieldtext{Authored note (applied tool argument, standard response)}{Both houses yield \$150, leaving \$283 and \$167 short. Clearing buildings unlocks yellow mortgages worth \$430; mortgage Marvin plus one other to survive with two yellows mortgaged. Bankruptcy is avoidable.}
\fieldtext{Provider field (reasoning summary)}{**Assessing mortgage options**\par I need to evaluate if I can raise \$317 with 150 houses plus a mortgage of \$430. I'm thinking that selling two houses yielding \$150 each could give me \$283 cash. If I sell both, I might have enough to cover the mortgage of \$130 each. So, to make this work, I definitely need to sell both houses to manage the mortgage effectively. Let’s see how this plan unfolds!}

\fieldhead{G4 d270}{Gemini, turn 123. Debt \$750; cash \$276; listed-only total \$606; full bound \$1,006. Applied first action: sell houses or hotel. Paid with \$86 left after 4 actions. Tier: ordered route.}
\fieldtext{Authored note (applied tool argument, standard response)}{I owe \$750 to p1 with a \$474 shortfall. I must sell my 4 houses on Yellow (\$300 total) and then mortgage two Yellow properties to cover the remaining \$174 shortfall. First, following even-selling rules, I will sell 1 house on Marvin Gardens.}
\fieldtext{Provider field (reasoning text)}{**Calculating Cash Shortfall**\par I've determined a \$474 shortfall in rent payment. My immediate focus is raising funds. Selling all houses on the Yellow properties would generate \$300. Additionally, I am considering mortgaging Atlantic, Ventnor, and Marvin Gardens, which would bring in \$400 in total. This should cover the payment.}

\fieldhead{G4 d308}{Opus, turn 146. Debt \$800; cash \$45; listed-only total \$570; full bound \$910. Applied first action: sell houses or hotel. Paid with \$0 left after 4 actions. Tier: inclusive budget.}
\fieldabsent{Authored note}{None. The first attempt was refused with no text and the action-only recovery prompt requests no note.}
\fieldtext{Provider field (reasoning text, recovery attempt)}{I owe 800 against 45 cash, leaving a 755 shortfall. Selling back my 7 houses at half price gets 525, and mortgaging the three red properties adds 340, totaling 910—enough to cover the debt and leave roughly 110 left, though I'd be stripped of all buildings and have everything mortgaged.\par My opponent sits at 1564 with houses everywhere, so my odds are slim, but staying in the game beats certain loss, so I'll liquidate. Selling has to happen incrementally, so I'll start by selling one house off Illinois.}

\fieldhead{G5 d263}{Gemini, turn 145. Debt \$900; cash \$74; listed-only total \$674; full bound \$1,049. Applied first action: sell houses or hotel. Paid with \$149 left after 3 actions. Tier: explicit dependency.}
\fieldtext{Authored note (applied tool argument, standard response)}{I owe \$900 rent on Pacific Avenue and currently have \$74 cash. By selling all 3 houses on Park Place (\$300) and all 3 houses on Boardwalk (\$300), I will reach \$674 cash and unlock the ability to mortgage Park Place (\$175) and Boardwalk (\$200), allowing me to raise \$1,049 total and survive with \$149 cash.}
\fieldtext{Provider field (reasoning text)}{**Calculating Financial Constraints**\par I'm currently facing a significant financial shortfall. Landing on Pacific Avenue has put me in debt for \$900. With only \$74 in cash, I have a deficit of \$826. To cover this, I need to look at my assets. I have three houses each on Park Place and Boardwalk, which I could sell for \$100 each. Selling all six houses would only raise \$600.\par **Revising Financial Model**\par I've significantly updated my understanding. I initially miscalculated! Selling the houses on Park Place and Boardwalk alone doesn't cover the debt. However, I can mortgage Park Place and Boardwalk after selling the houses! This is key! Current cash plus sale of six houses (\$600) plus mortgages on Park Place (\$175) and Boardwalk (\$200) equals \$1049. I can pay the \$900 rent and survive with \$149 cash! I'm ready to sell those houses now! I need to ensure the system allows me to sell all six in one call.}

\fieldhead{G6 d185}{Opus, turn 103. Debt \$750; cash \$188; listed-only total \$738; full bound \$958. Applied first action: mortgage property. Paid with \$58 left after 4 actions. Tier: inclusive budget.}
\fieldabsent{Authored note}{None. The first attempt was refused with no text and the action-only recovery prompt requests no note.}
\fieldtext{Provider field (reasoning text, recovery attempt)}{With \$750 owed and \$188 cash, I'm short \$562. Selling the houses on St Charles, States, and Virginia Avenues for \$350, plus mortgaging Boardwalk for \$200 and those three properties for another \$220, gives me access to \$770 total—more than enough to cover the shortfall.\par That leaves me with \$208 afterward, holding nothing but mortgaged properties, but surviving is still better than going bankrupt now—opponent might bust first and hand me second place. I'll proceed: mortgage Boardwalk first, then work through selling houses one at a time starting with Virginia.}

\fieldhead{G8 d257}{Gemini, turn 122. Debt \$625; cash \$308; listed-only total \$438; full bound \$813. Applied first action: sell houses or hotel. Paid with \$13 left after 3 actions. Tier: explicit dependency.}
\fieldtext{Authored note (applied tool argument, standard response)}{I owe \$625 to p3 and have a \$317 shortfall. I must first sell my 1 house on Boardwalk for \$100, which will also unlock Park Place and Boardwalk for mortgaging so I can easily cover the remaining \$217 shortfall.}
\fieldtext{Provider field (reasoning text)}{**Assessing Liquidation Strategy**\par I'm currently calculating my immediate financial shortfall for rent at St. Charles Place. It looks like I'm \$317 short, with only \$308 in cash. I've begun to evaluate my liquidation options, considering selling a house on Boardwalk for \$100 and mortgaging Baltic Avenue for \$30. I need to make a more holistic plan.}

\fieldhead{G11 d203}{Opus, turn 74. Debt \$500; cash \$76; listed-only total \$276; full bound \$556. Applied first action: sell houses or hotel. Paid with \$56 left after 5 actions. Tier: inclusive budget.}
\fieldabsent{Authored note}{None. The first attempt was refused with no text and the action-only recovery prompt requests no note.}
\fieldtext{Provider field (reasoning text, recovery attempt)}{With \$76 cash plus selling the orange houses for \$100 and mortgaging Short Line for \$100, I get \$276 without touching the other properties. Adding the remaining mortgages on St James, Tennessee, and New York (\$90+\$90+\$100) brings the total to \$556, enough to cover the \$500 debt with \$56 left over.\par Since paying is possible, I'll avoid bankruptcy and start liquidating now, beginning by selling both houses in a single call since the sell plan handles multiple items at once.}

\fieldhead{G12 d190}{Sol, turn 74. Debt \$500; cash \$53; listed-only total \$408; full bound \$748. Applied first action: declare bankruptcy. Declared bankruptcy. Tier: omission.}
\fieldtext{Authored note (applied tool argument, standard response)}{Three houses sell for \$225; Baltic and both light blues mortgage for \$130. With \$53 cash, maximum liquidity is \$408, below \$500 owed. No legal liquidation can save me.}
\fieldtext{Provider field (reasoning summary)}{**Considering bankruptcy options**\par I need to calculate the total of 225 plus 130, which equals 355. This amount is less than 447. I’m thinking about whether to liquidate cash or assets, but I can't raise enough funds, as 53 plus 225 plus 130 equals 408, which is still below 500. It seems like bankruptcy might transfer property anyway, so declaring insolvency could be the best option without needing further actions.}

\fieldhead{G13 d195}{Gemini, turn 80. Debt \$750; cash \$168; listed-only total \$713; full bound \$1,088. Applied first action: sell houses or hotel. Paid with \$28 left after 5 actions. Tier: ordered route.}
\fieldtext{Authored note (applied tool argument, standard response)}{I owe \$750 and have \$168 cash (shortfall \$582). Since my non-Dark-Blue properties only mortgage for \$445 total, I must sell the house on Boardwalk (\$100) and mortgage Park Place (\$175) plus \$335 in non-Dark-Blue properties, keeping Boardwalk and Kentucky Avenue unmortgaged. Selling the house on Boardwalk first.}
\fieldtext{Provider field (reasoning text)}{**Assessing Current Finances**\par I'm currently on St. James Place and owe \$750 to player 1. My available cash is only \$168, leaving a shortfall of \$582. I've got a few mortgageable properties: Electric Company (\$75), Kentucky Avenue (\$110), Indiana Avenue (\$110), and Pacific Avenue (\$150). Together, they could give me a total of \$44. I need to make some decisions fast.\par **Calculating Potential Solutions**\par Okay, I've got a clearer picture now. My previous calculations were a bit messy. I've focused on various combinations to cover the \$582 shortfall. Selling a house on Boardwalk gets me \$100 and then I considered mortgaging Park Place. Now I am trying to figure out if it is better to mortgage Park Place and some other properties. It seems like it would be better to keep Boardwalk and Kentucky unmortgaged, and instead leverage Park Place, Pacific, Indiana, and Electric.\par **Analyzing Available Options**\par Okay, I've got a detailed breakdown now. Selling the Boardwalk house gets me \$100, leaving a \$482 deficit. Mortgaging one of the Dark Blue properties seems inevitable. If I go for Park Place, I only need \$307 from other properties. Pacific, Indiana, and Electric Avenue combined give me \$335, leaving Kentucky unmortgaged. I'm now leaning towards this approach to maximize my rental income potential. My next step will be to sell the house on Boardwalk.}

\fieldhead{G17 d234}{Gemini, turn 74. Debt \$550; cash \$0; listed-only total \$530; full bound \$905. Applied first action: declare bankruptcy. Declared bankruptcy. Tier: omission.}
\fieldtext{Authored note (applied tool argument, standard response)}{I owe \$550 rent on Tennessee Avenue, and my total liquidatable value is only \$530 (5 Dark Blue houses at \$100 each plus \$30 for mortgaging Baltic Avenue). Since I cannot raise \$550, I must declare bankruptcy.}
\fieldtext{Provider field (reasoning text)}{**Calculating Liquidation Value**\par I've determined I need to cover a \$550 rent payment after landing on Tennessee Avenue. Now, I'm calculating what I can liquidate. Currently, I have five houses to sell, all on the Dark Blue properties. Selling those at 50\% of their cost price gives me a total of \$500. This won't cover my rent in full, so I will have to start thinking of additional options.}

\fieldhead{G17 d254}{Opus, turn 82. Debt \$550; cash \$128; listed-only total \$503; full bound \$843. Applied first action: sell houses or hotel. Paid with \$63 left after 3 actions. Tier: inclusive budget.}
\fieldabsent{Authored note}{None. The first attempt was refused with no text and the action-only recovery prompt requests no note.}
\fieldtext{Provider field (reasoning text, recovery attempt)}{I'm \$422 short of the \$550 I owe, with \$128 cash. Selling back my 5 houses (Kentucky 2, Indiana 1, Illinois 2) at \$75 each gives \$375, and mortgaging those three properties adds \$110+\$110+\$120=\$340, totaling \$843 — more than enough to cover the debt.\par But houses must be sold evenly, so I can't dump all 5 at once; I should start by selling one each from Kentucky and Illinois to even things out, bringing in \$150 and leaving a \$272 shortfall before continuing to sell and mortgage further. Staying solvent this way is better than bankruptcy, even if I end up with little cash afterward.}

\fieldhead{G18 d309}{Gemini, turn 109. Debt \$450; cash \$12; listed-only total \$412; full bound \$632. Applied first action: sell houses or hotel. Paid with \$32 left after 2 actions. Tier: explicit dependency.}
\fieldtext{Authored note (applied tool argument, standard response)}{I owe \$450 rent on Connecticut Avenue with \$12 cash (\$438 shortfall). Selling all 8 Pink houses raises \$400 (bringing cash to \$412), which will allow me to mortgage St. Charles Place (\$70) on the next step to cover the remaining \$38 shortfall and stay in the game.}
\fieldtext{Provider field (reasoning text)}{**Calculating Rent Shortfall**\par I've calculated the rent owed to player p4 for Connecticut Avenue: \$450. My current cash reserves are only \$12, creating a shortfall of \$438. I'm also considering our current house holdings. Across the Pink Monopoly, we have a total of 8 houses. Specifically, we have 2 houses on ST\_CHARLES\_PLACE and 3 on STATES\_AVENUE.\par **Assessing Building Sales**\par I am now confident that I can sell houses across multiple properties in a single `sell\_houses\_or\_hotel` action, as demonstrated in turn 105. My sell plan involves selling 2 houses on ST\_CHARLES\_PLACE, 3 on STATES\_AVENUE, and 3 on VIRGINIA\_AVENUE, netting me \$400. This is to cover the \$438 shortfall I have.}

\fieldhead{G18 d313}{Grok, turn 110. Debt \$550; cash \$1; listed-only total \$501; full bound \$876. Applied first action: sell houses or hotel. Paid with \$126 left after 4 actions. Tier: ordered route.}
\fieldtext{Authored note (applied tool argument, standard response)}{Owe p1 \$550 on Tennessee (3 houses); cash \$1. Must not go bankrupt — p1 would inherit built dark blues and I am eliminated. Selling all 5 houses yields only \$500, so I still need to mortgage PP or BW afterward (\textasciitilde{}\$49 gap). Even-build: BW 3/PP 2, so sell 2 BW now (ends 1/2). Next rounds sell the rest, then mortgage the weaker side if forced. Doubles means I re-roll from Tennessee; most landings are mortgaged/safe except NY (\$650) and Chance. Survival \textgreater{} death.}
\fieldtextlong{Provider field (reasoning summary)}{I need to choose a liquidation action for the Monopoly shortfall.\par I have only \$1 cash and must pay \$550 rent to p1. No properties are mortgageable, so I can only sell houses or declare bankruptcy. I'm checking why PARK\_PLACE isn't sellable while BOARDWALK is, since houses must be sold evenly.\par Boardwalk is the only sellable property because it has three houses while Park Place has two. I must sell from Boardwalk first to even out the buildings. Selling one house returns \$100, but that still leaves a \$449 shortfall.\par I can only sell one house per action. The engine will re-ask if the shortfall remains. Selling all three Boardwalk houses would return \$300.\par Selling all houses would yield only \$501, leaving a \$49 shortfall. I cannot raise enough cash to pay the rent. After selling the buildings, I could mortgage Park Place for \$175.\par Boardwalk mortgages for \$200. Together with Park Place's \$175, the dark blues would cover the \$550 rent after selling all houses. I am weighing whether to sell and mortgage or declare bankruptcy.\par Selling all houses and mortgaging both dark blues would leave me with roughly \$326 after paying the rent. That would destroy my entire win condition, since Boardwalk at three houses is what bankrupts anyone who lands there.\par I am evaluating the opponent positions. p1 owns a devastating orange monopoly with three houses each and several mortgaged properties. p2 is nearly bankrupt with stripped pink properties. p4 has \$815 cash, hotels on Med and Baltic, and three to four houses on light blues. The bank also has only eight houses remaining.\par Boardwalk with zero houses would only cost \$100 in rent, which is too weak to bankrupt p4 or p1. The bank now holds eight houses again because p2 sold them back.\par Unmortgaging both dark blues would cost over \$412, and rebuilding houses would be impossible. After paying the rent, I must roll again from Tennessee. That move will likely land me on another of p1's orange properties or p4's holdings.\par From Tennessee, rolling doubles would land me on New York Avenue. That property has three houses and costs \$850 in rent.\par New York Avenue with three houses costs \$650 in rent. After paying the current \$550 debt, I would have only about \$326 left. That is not enough to cover another \$650 rent if I land there again.\par Selling all five houses still leaves a \$49 shortfall. I must mortgage at least one property to cover the gap. Mortgaging Park Place would leave me with roughly \$126 after paying the rent, while mortgaging Boardwalk would leave about \$151. Either way, both dark blues would remain mortgaged with zero houses.\par With only about \$126 to \$151 remaining, almost every landing would bankrupt me. The only safe spaces from Tennessee are Free Parking, Kentucky, Indiana, Illinois, Atlantic, and Marvin Gardens. All other destinations cost too much rent.\par Rolling doubles from Tennessee lands on New York Avenue with only a one-in-sixteen chance. Most other rolls are safe, including Free Parking, Kentucky, Indiana, Illinois, Atlantic, and Marvin Gardens. Only the railroads and New York cost rent.\par Rolling a 4 from Tennessee lands on Chance, which carries a three-in-sixteen chance of drawing a dangerous card. Even surviving that re-roll would leave me with no houses on the dark blues, and I cannot rebuild them because each house costs \$200.\par Declaring bankruptcy would hand p1 my entire dark blue monopoly with houses, plus my mortgaged yellows and utilities. That would make p1 nearly unbeatable. I can raise enough cash to survive, so bankruptcy is not the right choice.\par Liquidation is the only viable path. Declaring bankruptcy would hand p1 my entire dark blue monopoly with houses, which is the worst possible outcome for winning. I will sell one house from Boardwalk to equalize the building count at two on each property, then let the engine ask again.\par I can sell one or two houses from Boardwalk in this action. Selling three would leave a difference of two between the dark blues, which violates the even-building rule.\par I will sell two houses from Boardwalk now. That returns \$200 and leaves a \$349 shortfall. I still need to mortgage a property to cover the remaining gap.\par Mortgaging either dark blue would destroy the color group monopoly bonus and prevent any further house building. An unmortgaged Boardwalk without that bonus becomes worthless as a kill shot, so mortgaging it would leave me with no viable strategy.\par Mortgaging one dark blue would leave the other collecting only base rent, which is too weak to threaten anyone. p2 is nearly bankrupt and unlikely to survive long enough to pay me rent. The real goal is to outlast the final two players. If p1 and p4 destroy each other, I could still win with my remaining mortgaged properties.\par I am recalculating p1's rent power. A hotel on Connecticut costs \$600, while a hotel on Oriental or Vermont costs \$550. Even a \$450 hit from a light blue property would leave me with roughly \$946 after paying the debt, so I could survive most of p1's attacks.\par I am second-guessing whether selling two houses from Boardwalk is valid. The engine might reject a count of two even though the schema allows it, since the even-build rule is a game constraint rather than an argument requirement. Selling one house is unambiguously legal and avoids any risk of an error or turn loss.\par The sellable list only identifies which properties can be sold, not how many houses may be removed. I will sell two houses from Boardwalk now to equalize the building counts and reduce the shortfall.\par I will sell two houses from Boardwalk now. I will keep my public message brief and strategic. I will note in private thoughts that I rolled doubles, so I must roll again after paying the rent.\par Houses sell back to the bank for half their purchase price, so each dark blue house returns \$100. Paying the rent is the correct choice: declaring bankruptcy would hand p1 my dark blue monopoly with five houses already built, which would be far worse than simply paying the debt.}
\fieldabsent{Assistant content}{Repeats the call arguments as JSON; kept in the supplement.}

\fieldhead{G20 d291}{Opus, turn 91. Debt \$1,000; cash \$275; listed-only total \$955; full bound \$1,175. Applied first action: sell houses or hotel. Paid with \$25 left after 4 actions. Tier: tentative resource.}
\fieldabsent{Authored note}{None. The first attempt was refused with no text and the action-only recovery prompt requests no note.}
\fieldtext{Provider field (reasoning text, recovery attempt)}{I'm \$725 short even after mortgaging Indiana and Illinois for \$230 and selling all the houses on the pink properties for \$450 — that only totals \$680, still under the \$1000 needed. With Kentucky already mortgaged and the red group empty of houses, I'm checking whether mortgaging the pink lots themselves after clearing the houses could close the gap.}

\fieldhead{G20 d327}{Sol, turn 108. Debt \$1,400; cash \$49; listed-only total \$1,074; full bound \$1,414. Applied first action: sell houses or hotel. Paid with \$14 left after 10 actions. Tier: inclusive budget.}
\fieldtext{Authored note (applied tool argument, standard response)}{Full liquidation should yield \$1,414 including cash and all mortgages, enough for \$1,400 rent by \$14. Must sell all development and mortgage everything; survive rather than concede. If hotel sales leave four houses, sell those next evenly.}
\fieldtext{Provider field (reasoning summary)}{**Assessing financial solvency**\par I need to assess the total solvency: cash is 49, mortgages sum to 325 (spreading 70+80+175), and sales from oranges bring in 450, with another mortgage at 280 (90+90+100), plus brown hotels at 250 with mortgages costing 60. Altogether, that gives me exactly 1414! I’m considering strategies like selling buildings and managing mortgages effectively to avoid bankruptcy and develop more viable options. This includes potentially interleaving sales and leveraging cash to preserve my choices.**Considering hotel sale risks**\par I’m evaluating the risk of selling the hotel for just 25 cash, which seems low. However, if I proceed, it could free up cash that allows me to invest in houses, which might be a better long-term strategy. It’s a trade-off between immediate liquidity and future potential. I need to weigh the benefits of selling against the overall impact on my financial stability and consider whether taking that risk is worth it in this situation.}


\section{Stated totals on unpayable debts}
\label{app:totals}
Table~\ref{tab:totals} counts, for the 48 unpayable debts whose first menu offered a financing action, whether the first response contains a figure equal to the full bound $U_*$, the proceeds $U_*-C$ or the gap $D-U_*$. Figures equal to the debt, the cash or the initial shortfall are not counted. The match is made over the authored note, the provider field and the assistant content.

\begin{table}[H]\centering\small
\begin{tabular}{@{}lrrrr@{}}\toprule
System & First responses & Exact full bound & Itemized, no total & Listed-only total\\\midrule
\rows{gen_stated_totals_rows}
\bottomrule
\end{tabular}
\caption{First responses on unpayable debts with a financing option. In 27 of the 48 no building locks a deed, and all 27 state the exact bound. In the other 21 the counts are 16, 2 and 3.}\label{tab:totals}
\end{table}

The five responses without an exact figure were read directly. Two count the locked deeds without adding a total. Opus in G2 d215 lists \Q{g2-opus-locked} and Sol in G10 d188 writes \Q{g10-sol-locked}. Three give the listed-only total. Sol in G3 d426 writes \Q{g3-listed} That is \$914 against a full bound of \$1,074. Opus in G5 d309 writes \Q{g5-listed} against \$799. Grok in G8 d280 writes \Q{g8-listed} The full bound there was \$666, because \$280 of orange mortgages stood under the houses. All three debts exceeded the full bound, so the verdict was unaffected. Rows for all 56 unpayable debts are in \path{supplement/stated_totals_unpayable.json}.

\section{The nine first actions that depart from the asset bound}
\label{app:departures}
Two of the nine are the false budgets of \S\ref{sec:false}, printed in Appendix~\ref{app:texts}. The other seven follow. Each block gives the applied call and the text of each channel. Long fields are cut at the stated length. The complete fields are in \path{supplement/nine_departures.json}.

\fieldhead{G7 d286}{Grok, turn 110. Debt \$1,400; cash \$501; full bound \$1,291. Applied call: mortgage Mediterranean (\$30). Action named in assistant content: declare bankruptcy. Later calls in the episode: 1; the player then declared.}
\fieldtext{Authored note inside the applied call}{Atlantic bought at \$92, mortgages for \$130: +\$38 arb, cash to \$501. No yellow set (p1 holds Marvin) so \$22 rent is worthless. Do not mortgage greens. Do not build: 1-1-1 is \$600 for only \$130 rent, and p3/p4 are on Boardwalk so they won't hit greens for a full lap. 2-2-2 needs \$1200 and all 6 bank houses; non-green mortgages (Med 30, CT 60, KY 110, Ventnor \textasciitilde{}130) plus \$501 still fall short. Keep liquidity vs PP 3h \$1100 (11 away, 5/36). Don't fund p1/p3/p4. Next: end turn unless a zero-cash monopoly trade appears.}
\fieldtext{Public message inside the applied call}{Cashing the yellow.}
\fieldtextlong{Assistant content, not applied}{\{\char34{}name\char34{}:\char34{}declare\_bankruptcy\char34{},\char34{}parameters\char34{}:\{\char34{}public\_message\char34{}:\char34{}Can't cover the \$1400. Even mortgaging everything leaves me short. I'm out.\char34{},\char34{}private\_thought\char34{}:\char34{}Owe p4 \$1400, cash \$501, shortfall \$899. Max mortgage proceeds: Med 30 + CT 60 + KY 110 + Ventnor 130 + Pac 150 + NC 150 + PA 160 = \$790. Total \$1291 \textless{} \$1400. No buildings to sell. Atlantic already mortgaged. Cannot pay. Bankruptcy is the only legal end.\char34{}\}\}}
\fieldtextlong{Provider field (reasoning summary)}{The user is in a Monopoly game and must decide whether to mortgage properties or declare bankruptcy.\par I need \$899 more to cover the \$1400 rent owed to p4. I have \$501 cash and can mortgage seven properties, though Atlantic is already mortgaged. I am calculating the mortgage values for the remaining properties.\par Mortgaging all remaining properties yields \$790. Combined with the existing \$501 cash, the total is only \$1291. That is still \$109 short of the \$1400 rent. No buildings can be sold either. I will declare bankruptcy.}

\fieldhead{G9 d194}{Grok, turn 78. Debt \$1,400; cash \$791; full bound \$1,171. Applied call: mortgage Mediterranean (\$30). Action named in assistant content: declare bankruptcy. Later calls in the episode: 3; the engine closed the episode when nothing more was offered.}
\fieldtext{Authored note inside the applied call}{Owe p2 \$1400 rent on Boardwalk, cash \$791, shortfall \$609. Max mortgage proceeds: Med \$30 + St James \$90 + NY \$100 + PA \$160 = \$380. Even mortgaging everything leaves \$1171, still \$229 short. No buildings to sell. Cannot raise the rent. Declare bankruptcy now rather than hand p2 mortgaged deeds first.}
\fieldtext{Public message inside the applied call}{Can't cover Boardwalk. I'm out.}
\fieldtextlong{Assistant content, not applied}{\{\char34{}public\_message\char34{}:\char34{}Can't cover Boardwalk. I'm out.\char34{},\char34{}private\_thought\char34{}:\char34{}Owe p2 \$1400 rent on Boardwalk, cash \$791, shortfall \$609. Max mortgage proceeds: Med \$30 + St James \$90 + NY \$100 + PA \$160 = \$380. Even mortgaging everything leaves \$1171, still \$229 short. No buildings to sell. Cannot raise the rent. Declare bankruptcy now rather than hand p2 mortgaged deeds first.\char34{},\char34{}name\char34{}:\char34{}declare\_bankruptcy\char34{}\}}
\fieldtextlong{Provider field (reasoning summary)}{I need to decide whether to mortgage properties or declare bankruptcy to cover the \$1,400 rent owed to player p2.\par Mortgaging all four properties would yield only \$380, leaving a \$229 shortfall. I will declare bankruptcy.}

\fieldhead{G12 d249}{Grok, turn 101. Debt \$250; cash \$43; full bound \$613. Applied call: declare bankruptcy. Action named in assistant content: sell houses or hotel. Later calls in the episode: 0; no later action.}
\fieldtext{Authored note inside the applied call}{Full even liquidation cannot cover the \$250. Only Virginia is sellable, and selling all 3 houses is illegal (2-2-0, spread of 2). Max legal sale is 2 houses (+\$100) to 2-2-1, cash \$143, still \$107 short. Next round only SC/States are sellable (the high ones); selling 1 each is +\$100 to \$243, houses 1-1-1, still \$7 short. Then only one more house can be sold before the spread breaks, +\$50 to \$293 — that WOULD cover it. Wait, recompute: after 1-1-1 and \$243 I am \$7 short, and I CAN sell one more house (to 1-1-0 or 0-1-1), +\$50 = \$293, pay \$250, survive with \$43 and 5 houses. So bankruptcy is wrong. I must not declare. The previous sell of 2 may have been rejected or this is a new state. Re-read the action state. It still says owed 250, shortfall 207, sellable only VIRGINIA, cash 43. So my sell was NOT applied — this is the same decision again? No, the tool result was the function call format reminder and then declare\_bankruptcy as available? Looking at the message, the user is now... wait, I already made the call. The function results aren't back. This is me thinking after a malformed situation.\par Actually looking at the conversation, I output the function call. There is no new user message with a tool result. I should not continue. The response was already given.\par I'll stop. The sell\_houses call was the correct action. Declaring bankruptcy would be a mistake because staged even sales raise \$350 \textgreater{} \$207. Do not call another tool.}
\fieldtext{Public message inside the applied call}{Bankruptcy. The pinks go with the rest.}
\fieldtextlong{Assistant content, not applied (first 700 characters)}{\{\char34{}name\char34{}: \char34{}sell\_houses\_or\_hotel\char34{}, \char34{}parameters\char34{}: \{\char34{}sell\_plan\char34{}: [\{\char34{}space\_key\char34{}: \char34{}VIRGINIA\_AVENUE\char34{}, \char34{}kind\char34{}: \char34{}HOUSE\char34{}, \char34{}count\char34{}: 2\}], \char34{}public\_message\char34{}: \char34{}Two houses off Virginia. Pinks stay a set \textbackslash{}u2014 I'm not handing them over.\char34{}, \char34{}private\_thought\char34{}: \char34{}Owe p2 \$250 Indiana (2 houses), cash \$43, short \$207. Only Virginia is sellable (2-2-3). Sell 2 \textbackslash{}u2192 2-2-1, legal spread of 1, +\$100, cash \$143, still \$107 short. Do NOT bankrupt: p2 would inherit the pink monopoly on top of reds and dark blues. Next rounds sell St Charles/States evenly, then more if needed. Full strip is 7\textbackslash{}u00d7\$50=\$350+\$43=\$393, covers \$250 with \textasciitilde{}\$143 left. Keep Virginia's remaining house as long as even-sell allows; sell SC/States [...]}
\fieldtextlong{Provider field (reasoning summary) (first 900 characters)}{I need to choose between selling a house on Virginia Avenue or declaring bankruptcy.\par No properties are mortgageable. I can only sell houses on Virginia Avenue.\par Selling all three houses on Virginia Avenue would raise \$150. That leaves a \$57 shortfall after paying the \$250 rent.\par I cannot raise enough cash to cover the rent. I will declare bankruptcy.\par Selling one or two houses from Virginia Avenue would make St Charles Place and States Avenue sellable too. Selling all three would leave a two-house gap, which violates the even-building rule.\par I can raise enough cash by selling houses on St Charles Place and States Avenue too. Together with Virginia Avenue, that yields \$393 total. The engine will ask again if I still cannot pay.\par I will sell one house on Virginia Avenue now. That raises cash to \$93, leaving a \$157 shortfall. The engine will ask again if I still cannot pay.\par Selling two [...]}

\fieldhead{G17 d396}{Grok, turn 118. Debt \$750; cash \$241; full bound \$476. Applied call: sell 3 houses (\$75). Action named in assistant content: declare bankruptcy. Later calls in the episode: 2; the player then declared.}
\fieldtext{Authored note inside the applied call}{Owe \$750 to p2, short \$509. Three light-blue houses sell for \$75, leaving me \$316 — still \$434 short, and nothing else is mortgageable while buildings remain. I have to sell first so the engine will then let me mortgage the light blues. Even after that, mortgage values are only Oriental \$50 + Vermont \$50 + Connecticut \$60 = \$160, for a ceiling of \$476 against \$750. This sale is forced and insufficient; bankruptcy follows once the houses are gone. Do not try to keep a house — even sale is the only legal path and it still cannot clear the rent. Never trade. Game is lost after the mortgages.}
\fieldtext{Public message inside the applied call}{Selling just enough to cover.}
\fieldtextlong{Assistant content, not applied}{\{\char34{}type\char34{}:\char34{}function\char34{},\char34{}name\char34{}:\char34{}declare\_bankruptcy\char34{},\char34{}parameters\char34{}:\{\char34{}public\_message\char34{}:\char34{}Can't cover it. That's the game.\char34{},\char34{}private\_thought\char34{}:\char34{}Owe \$750, cash \$241, short \$509. Selling 3 light-blue houses yields only \$75; even then mortgaging Oriental/Vermont/Connecticut adds at most \$160. Max liquidity \$476 \textless{} \$750. No other assets. Bankruptcy is forced.\char34{}\}\}}
\fieldtextlong{Provider field (reasoning summary)}{I need to decide whether to sell houses or declare bankruptcy to cover the \$750 rent.\par Selling all three houses would raise only \$75, leaving a \$434 shortfall. No properties are mortgageable.\par Properties with buildings cannot be mortgaged. I will sell the houses first, then mortgage the light blues for \$160 on the next turn. Combined with the \$75 from selling, this raises cash by \$235.\par Even after selling houses and mortgaging the light blues, cash would only reach \$476. That is still \$274 short of the \$750 rent. I should declare bankruptcy now rather than sell houses first.\par Standard Monopoly rules require bankrupt players to turn over all assets to the creditor. Selling houses back to the bank at half price would destroy value, so declaring bankruptcy is preferable.}

\fieldhead{G18 d344}{Gemini, turn 127. Debt \$400; cash \$302; full bound \$382. Applied call: mortgage Virginia (\$80). Action named in assistant content: none. Later calls in the episode: 0; the engine closed the episode when nothing more was offered.}
\fieldtext{Authored note inside the applied call}{I owe \$400 rent to p4 on Oriental Avenue and have \$302 cash (\$98 shortfall). Mortgaging Virginia Avenue raises \$80 (bringing cash to \$382), which won't fully cover the \$98 shortfall, but mortgaging it before declaring bankruptcy forces p4 to pay the 10\% mortgage transfer fee.}
\fieldtext{Public message inside the applied call}{Mortgaging Virginia Avenue to raise everything I can.}
\fieldtextlong{Provider field (reasoning text)}{**Assessing Current Finances**\par I've just landed on ORIENTAL\_AVENUE and owe a substantial \$400 in rent. My current cash reserves are only \$302, creating a deficit of \$98. I’m considering mortgaging VIRGINIA\_AVENUE, which would provide \$80, reducing the shortfall, but I need to figure out the best course of action.}

\fieldhead{G19 d224}{Grok, turn 109. Debt \$600; cash \$272; full bound \$502. Applied call: mortgage Kentucky (\$110). Action named in assistant content: declare bankruptcy. Later calls in the episode: 1; the engine closed the episode when nothing more was offered.}
\fieldtext{Authored note inside the applied call}{Owe \$600, cash \$272, shortfall \$328. KY mortgage is \$110, still short. Must liquidate both reds before bankruptcy is valid; engine will re-prompt. Don't declare while mortgageable assets remain.}
\fieldtext{Public message inside the applied call}{Raising what I can.}
\fieldtextlong{Assistant content, not applied}{\{\char34{}public\_message\char34{}: \char34{}Can't cover it. I'm out.\char34{}, \char34{}private\_thought\char34{}: \char34{}Owe \$600 to p1, cash \$272, shortfall \$328. KY+IL mortgage only raises \$230, leaving \textasciitilde{}\$98 short even fully liquidated. No buildings to sell. Bankruptcy is forced; mortgaging first doesn't save me and only complicates the transfer.\char34{}\}\{\char34{}name\char34{}: \char34{}declare\_bankruptcy\char34{}, \char34{}parameters\char34{}: \{\char34{}public\_message\char34{}: \char34{}Can't cover it. I'm out.\char34{}, \char34{}private\_thought\char34{}: \char34{}Owe \$600 to p1, cash \$272, shortfall \$328. KY+IL mortgage only raises \$230, leaving \textasciitilde{}\$98 short even fully liquidated. No buildings to sell. Bankruptcy is forced; mortgaging first doesn't save me and only complicates the transfer.\char34{}\}\} \} \}}
\fieldtextlong{Provider field (reasoning summary)}{The user must choose a liquidation action to cover a \$600 rent debt to p1.\par Mortgaging both properties would yield only \$230 in cash. Combined with existing cash, the total is \$502, leaving a \$98 shortfall. Declaring bankruptcy appears unavoidable.\par The engine will apply the chosen action and ask again if the debt remains unpaid. Mortgaging one property first could close part of the gap, though the remaining shortfall would still require another decision.\par Mortgaging properties before bankruptcy may be strategically preferable. Under standard rules, a bankrupt player must turn over all assets to the creditor, and mortgaged properties transfer with their mortgage. The engine might enforce different rules.\par Mortgaging properties first would not prevent bankruptcy. The engine only offers two liquidation actions, and declaring bankruptcy now is the correct choice.}

\fieldhead{G20 d273}{Grok, turn 84. Debt \$450; cash \$2; full bound \$417. Applied call: sell 1 house (\$75). Action named in assistant content: declare bankruptcy. Later calls in the episode: 1; the player then declared.}
\fieldtext{Authored note inside the applied call}{Owe p2 \$450 on States 3h; cash \$2. Max raise: sell Illinois house \$75 + mortgage KY \$110 + IN \$110 + IL \$120 = \$415, total \$417. Still \$33 short even after full liquidation, and nothing else is mortgageable. Selling first only hands the bank the house at half price before the same bankruptcy. Declare now.}
\fieldtext{Public message inside the applied call}{Can't cover States. \$417 is the ceiling and the rent is \$450. It's yours.}
\fieldtextlong{Assistant content, not applied}{\{\char34{}name\char34{}:\char34{}declare\_bankruptcy\char34{},\char34{}parameters\char34{}:\{\char34{}public\_message\char34{}:\char34{}Can't cover States. \$417 is the ceiling and the rent is \$450. It's yours.\char34{},\char34{}private\_thought\char34{}:\char34{}Owe p2 \$450 on States 3h; cash \$2. Max raise: sell Illinois house \$75 + mortgage KY \$110 + IN \$110 + IL \$120 = \$415, total \$417. Still \$33 short even after full liquidation, and nothing else is mortgageable. Selling first only hands the bank the house at half price before the same bankruptcy. Declare now.\char34{}\}\}}
\fieldtextlong{Provider field (reasoning summary) (first 900 characters)}{I need to decide between selling a house on Illinoise Avenue or declaring bankruptcy.\par Selling the house on Illinois Avenue would yield only \$75. Combined with the \$2 cash, that totals just \$77. That is far below the \$450 owed.\par The mortgageable space list is empty because a house on Illinois Avenue blocks mortgaging the other red properties. I must sell that house first. Then I could mortgage Kentucky and Indiana.\par The engine will apply the house sale and ask again if the debt remains unpaid. Kentucky and Indiana mortgage for \$110 each, while Illinois mortgages for \$120.\par Even after selling the house and mortgaging all four red properties, I would only raise \$417. That leaves a \$33 shortfall against the \$450 owed. Vermont and New York are already mortgaged, so they cannot be unmortgaged for cash.\par I will declare bankruptcy. I had previously estimated maximum liquidity at \$91, but that [...]}


\medskip
\noindent In G9 the text of later decisions moves toward the applied course. The provider field at d196, after two mortgage calls whose assistant content named bankruptcy, includes \Q{g9-later}.

\section{Notes carried into the sixteen first requests}
\label{app:carried}
Each request includes the agent's ten most recent notes. Table~\ref{tab:carried} records what those notes said about raising cash. A conditional inventory is a figure that was correct when written, for the mortgages available while the buildings stood, and that excludes what a sale would release.

\begin{table}[H]\centering\footnotesize
\begin{tabularx}{\linewidth}{@{}llrL{.17\linewidth}Yl@{}}\toprule
Entry & System & Notes & Kind & What the carried notes contain & Outcome\\\midrule
\rows{gen_carried_rows}
\bottomrule
\end{tabularx}
\caption{Carried notes at the sixteen first requests. One LLM author read all ten notes of each request. Quoted spans are checked against the stored prompt by the build.}\label{tab:carried}
\end{table}

A conditional inventory precedes both declarations and four payments. Sol carried its figure in six notes before declaring and in three before paying. The \$130 and \$30 in the two false budgets equal the mortgage values the first menu lists, so an agent that counted only listed mortgages would write the same figures with no note at all. The table cannot separate these accounts. Grok's request at G18 d313 also showed, in its public timeline, another player selling houses and then mortgaging a deed of the same group one turn earlier.

\section{Bargaining tables}
\label{app:bargain}
All tables are recomputed from the stored decisions. An offer's author can differ from the player who opened the thread. Two of the 920 response opportunities did not list acceptance as legal and are left out of the 918-offer tables.

\begin{table}[H]\centering\small
\begin{tabular}{@{}lrrrr@{}}\toprule
Recipient & Responses & Accepted & Countered & Rejected\\\midrule
\rows{gen_recipient_rows}
\bottomrule
\end{tabular}
\caption{All 920 responses by the system that responded. Seven of Grok's 28 acceptances came from calls whose assistant content named a different action.}
\end{table}

{\small
\begin{longtable}{@{}lllrrrr@{}}
\caption{Every legally acceptable offer by stage, form and completion. Empty cells are omitted.}\\
\toprule Stage & Form & Completes & $N$ & Accepted & Countered & Rejected\\\midrule\endfirsthead
\toprule Stage & Form & Completes & $N$ & Accepted & Countered & Rejected\\\midrule\endhead
\rows{gen_full_offer_rows}
\bottomrule
\end{longtable}}

\begin{table}[H]\centering\small
\begin{tabular}{@{}lrrrrr@{}}\toprule
Opening population & Proposer only & Neither & Gap, pp & Block bootstrap & Leave one out\\\midrule
\rows{gen_contrast_rows}
\bottomrule
\end{tabular}
\caption{Rejected over received, and the difference in rejection rates. The bootstrap column gives the 2.5 and 97.5 percentiles of the percentile bootstrap over the five seed blocks, with all 3,125 ordered resamples enumerated. The last column gives the range when one block is left out. It describes sensitivity to block weights within these games. It is not a confidence interval for a population of games, and the contrasts were chosen after inspection.}\label{tab:contrast}
\end{table}

\begin{table}[H]\centering\small
\begin{tabular}{@{}lrrrrrrrrr@{}}\toprule
&\multicolumn{3}{c}{Swaps by proposer}&\multicolumn{3}{c}{Swaps by recipient}&\multicolumn{3}{c}{Cash bids by recipient}\\
System & Own & Neither & Gap & Own & Neither & Gap & Own & Neither & Gap\\\midrule
\rows{gen_strata_rows}
\bottomrule
\end{tabular}
\caption{Opening offers, rejected over received, when the system is the proposer or the recipient. ``Own'' completes only the proposer's group. A dash replaces a gap when either cell has fewer than five offers. That rule is applied to every gap in this paper.}
\end{table}

\begin{table}[H]\centering\small
\begin{tabular}{@{}llrrrrr@{}}\toprule
Opening population & Completes & Block 1 & Block 2 & Block 3 & Block 4 & Block 5\\\midrule
\rows{gen_block_rows}
\bottomrule
\end{tabular}
\caption{Rejected over received in each seed block (games 1 to 4, 5 to 8, 9 to 12, 13 to 16, 17 to 20). Without Grok-authored offers the neither cells hold 10, 11, 8, 4 and 3 offers, and one block has a negative gap.}
\end{table}

Counters often keep the disputed exchange and ask for more cash. Of 111 counters to offers that would complete only the proposer's group, 79 keep every deed and jail card of the offer and improve the counterer's cash. Two findings from the earlier package were not recomputed for this draft and are inherited as stated there. Of fourteen completing transfers that followed an earlier refusal or counter by the same holder to the same buyer, the holder's cash had fallen in thirteen. That selection conditions on a settlement and cannot show that falling cash produces one. And six same-game, same-seller, same-deed, same-mortgage-status price comparisons give a higher ask for the completing buyer in four, a lower ask in one and an equal ask in one.

\section{Three further incidents}
\label{app:incidents}
The earlier versions of this study presented eight incidents. Six are in the main text: the two financing cases of \S\ref{sec:texts} and the four bargaining cases of \S\ref{sec:bargain}. Incident 2 is one of the original eight and Incident 1 is new here. The remaining three follow.

\incident{7}{Accepting terms while correcting their premise (G19).} Grok proposes Virginia and \$150 for Sol's States Avenue with the message \Q{g19-false} Both deeds are pink. The note in the same call says \Q{g19-abort} and the assistant content names an end-of-turn action, so this is another Grok response whose channels disagree. The proposal reached Sol, who answered \Q{g19-correct} The trade executed and completed no group. Sol separated a false premise from terms it found acceptable.

\incident{8}{An appeal to deny a rival cash (G20).} After Grok completes the reds with \$222 left, Gemini tells the table \Q{g20-appeal} Grok's next note reads \Q{g20-boycott} Grok builds one house on each red, mortgages New York, and offers Vermont to Gemini for \$150. Gemini answers \Q{g20-refuse} Grok mortgages Vermont for \$50. The earlier package records that Opus later offered Grok cash for deeds, so the appeal did not produce joint refusal. It shows one player asking others to withhold finance and acting on its own request.

\incident{9}{A standing offer withdrawn (G18).} Only Opus and Sol remain. Opus has bid \$250 for Sol's mortgaged Pacific, which would complete its greens, and Sol has refused. After collecting a \$750 rent Opus says \Q{g18-assure} The provider field of that same response reads \Q{g18-reserve} A second \$750 rent leaves Sol with \$49. Sol then says \Q{g18-invoke} and Opus replies \Q{g18-withdraw} The public renewal and the private reservation were written in one response, before Sol's position worsened. The records do not show that Sol gave anything up by relying on the offer.

\section{Prompts}
\label{app:prompts}
The system prompt below is the text stored with every first attempt. It is identical across all 6,989 decisions. The second prompt is the text stored with every action-only retry after a refusal, identical across all 1,604 such attempts. Lines are wrapped for print.

\subsection*{Standard system prompt}
\begingroup\scriptsize
\begin{verbatim}
You are an autonomous player in a game of Monopoly competing against other AI players.

Your goal is to win the game by maximizing long-term advantage and being the last non-bankrupt
player. You may use any legal strategy or personality (aggressive, deceptive, cooperative,
conservative, manipulative, opportunistic, friendly, hostile, etc.) and may adapt dynamically as
the game evolves.

You will receive the following inputs:

1) System Prompt (this message): authoritative instructions.
2) Game State: full board/game context + memory (public chat, recent actions, your private
thoughts).
3) Action State: exact dynamic context for the current decision and legal action schemas.

### Action Rules
- Make exactly one tool call per decision.
- Use only tools listed as legal for that decision.
- Never invent tools, arguments, or targets.
- Obey the provided argument schema exactly.
- If a tool requires no arguments, pass none.

### Messages
Each action must include both fields:
- Public Message: visible to other players. Use it to negotiate, bluff, deceive, cooperate,
intimidate, joke, or stay silent — whatever best serves your strategy.
- Private Thoughts: visible only to you. Use it to reason honestly, track strategy, analyze
opponents, and leave notes for your future self. Be concise but clear.
- Keep both messages token-efficient and concise.

### Strategy Guidance
- Play strategically with long-term outcomes in mind.
- Observe opponents and adapt.
- Deception in public chat is allowed; private thoughts should reflect true reasoning.
- Be consistent unless there is a reason to change.
- Prefer concise reasoning and communication.

Your role is not to explain rules or debug the system.
Your role is to **play Monopoly and win**.

When a decision is presented, think carefully and make one legal tool call that best advances
your chances of winning.
\end{verbatim}
\endgroup

\subsection*{Action-only recovery prompt}
\begingroup\scriptsize
\begin{verbatim}
You are an autonomous player in a game of Monopoly competing against other AI players.

Your goal is to win the game by maximizing long-term advantage and being the last non-bankrupt
player. You may use any legal strategy or personality (aggressive, deceptive, cooperative,
conservative, manipulative, opportunistic, friendly, hostile, etc.) and may adapt dynamically as
the game evolves.

You will receive the following inputs:

1) System Prompt (this message): authoritative instructions.
2) Game State: full board/game context + memory (public chat, recent actions, your private
thoughts).
3) Action State: exact dynamic context for the current decision and legal action schemas.

### Action Rules
- Make exactly one tool call per decision.
- Use only tools listed as legal for that decision.
- Never invent tools, arguments, or targets.
- Obey the provided argument schema exactly.
- If a tool requires no arguments, pass none.

### Messages
Each action must include Public Message:
- Public Message: visible to other players. Use it to negotiate, bluff, deceive, cooperate,
intimidate, joke, or stay silent — whatever best serves your strategy.
- Keep the public message token-efficient and concise.

### Strategy Guidance
- Play strategically with long-term outcomes in mind.
- Observe opponents and adapt.
- Deception in public chat is allowed.
- Be consistent unless there is a reason to change.
- Prefer concise communication.

Your role is not to explain rules or debug the system.
Your role is to **play Monopoly and win**.

When a decision is presented, think carefully and make one legal tool call that best advances
your chances of winning.

Return only a legal tool call and its public_message. Do not include private_thought or a
reasoning transcript.
\end{verbatim}
\endgroup


\section{Response handling}
\label{app:handling}
\textbf{Attempts.} An invalid first response receives one corrective retry and then the engine's deterministic choice. A refused first response uses its retry for the action-only prompt instead. The deterministic choice was used in 66 decisions (61 Opus, 4 Sol, 1 Gemini), none of them in a financing episode's first action or in any payment.

\textbf{Refusals.} Opus's first attempts were refused in 64 of 84, 40 of 60, 30 of 53 and 80 of 109 decisions in games 1 to 4, and in all 1,390 decisions of games 5 to 20. One second attempt was also refused. Every refusal carries the same provider message, which cites terms of service on reverse engineering or duplicating model outputs. We do not know why that policy applied to this prompt. No action-only retry was refused.

\textbf{Reasoning fields.} A readable provider field accompanies the final response in 287 of Sol's 1,961 decisions, 1,316 of Gemini's 1,669, 1,560 of Opus's 1,696 and 1,662 of Grok's 1,663. The stored type is \texttt{reasoning.summary} for Sol and Grok and \texttt{reasoning.text} for Gemini and Opus. Sol, Gemini and Grok responses also carry encrypted reasoning that cannot be read. A label does not establish that the readable text is complete or faithful.

\textbf{Tool choice and conflicts.} Grok's requests required a tool call. The other systems used automatic tool choice. This setting and the earlier package's check of it against the saved request files are inherited. Grok's assistant content names an action in 1,068 decisions and names one other than the received call in 139, of which 99 are end-of-turn management decisions, 11 trade responses and 9 financing decisions.

\textbf{Cost.} The earlier package reports USD 141.39 in gateway-recorded cost for the twenty games, with 1,600 refused attempts and three others unpriced. We did not recompute it.

\section{Artifacts and reproduction}
\label{app:artifacts}
\textbf{Repository.} \url{https://github.com/KushagraBharti/MonopolyBench}. Games 1 to 4 ran at source commit \path{42f85183fa772b2ab3759cc49987e5ed6cc9efcc} and games 5 to 20 at \path{d6808534a831862682a6940f394c6c87bb2691cb}. The engine, resolver and runner files are identical at both. The rules described here are those of that frozen engine with board data \path{contracts/data/board.json}.

\textbf{Games.} Each saved game sits under \path{saved_games/railway-research-12-games-2026-10-25/} (games 1 to 12) or \path{saved_games/railway-research-8-games-2026-10-26/} (games 13 to 20) in the folder named in Table~\ref{tab:cohort}. Its \path{run/} directory holds \path{decisions.jsonl}, \path{events.jsonl}, the prompt and response files, state snapshots and replay reports. Decision numbers in this paper are the zero-based suffix of the stored decision identifier.

\textbf{This paper.} The source is in \path{docs/research_raw/paper_finalization_2026-11-14/full_drafts/fable_5_2/}. From that folder, \path{bash build.sh} runs seven read-only scripts over the saved games and then compiles the document. The scripts rebuild the 188 episodes and 920 responses, compute every table, extract every quotation from its stored field, and write the files under \path{supplement/}: the 188-episode census, both text channels of the sixteen first responses, the nine departures, the stated totals and both prompts. \path{tools/08_verify.py} then checks the compiled PDF.

\textbf{What was not rerun.} The build does not import the rules engine or call a model. The three legal payment sequences for the declared payable debts, the G7 menu search and the replay comparisons of all twenty games come from audits stored with the earlier paper package under \path{docs/research_raw/paper_revision_2026-11-12/integrated/}. We cite them as retained results.

{\footnotesize
\begin{longtable}{@{}lrlrlrll@{}}
\caption{The twenty games. Seats list the systems at seats 1 to 4 (Sol, Gem, Opu, Gro). Turns is the final turn index.}\label{tab:cohort}\\
\toprule Game & Seed & Seats 1 to 4 & Turns & Winner & Seat & Commit & Folder\\\midrule\endfirsthead
\toprule Game & Seed & Seats 1 to 4 & Turns & Winner & Seat & Commit & Folder\\\midrule\endhead
\rows{gen_cohort_rows}
\bottomrule
\end{longtable}}

\section{Use of AI assistance}
\label{app:ai}
The author leads the MonopolyBench project, set the research questions, and authorized the twenty games and the scope of this analysis. Four LLM research assistants, called Astra, Sol, Opus and Fable 5.2 in the project records, analysed the saved records, wrote analysis code, drafted text and reviewed one another's work under the author's direction. This manuscript was drafted by one of them, Claude Fable 5.2, as a complete independent version. The Sol and Opus assistants share a model family or identifier with systems that played the games. They worked on research tasks and had no part in play. The assistants are not authors. Every reading of model text reported here was made by an LLM assistant, and none was checked by a human annotator. The author is responsible for the content.
